% Definition file for row def_3_4 -- "Walks".
% Auto-generated stub by scaffold/solve_chapter.py at row start. The
% manager agent (or a worker dispatched by it) fills the body below with the
% Lean-faithful definition statement(s). Stays close to the lecture notes.
%
% This file is a sibling of `main.tex` inside `<subsection>/tex/`. The
% `subfiles` package lets it compile both standalone (resolving its
% preamble via `main.tex`) and as part of `main.tex` via `\subfile{...}`.

\documentclass[main]{subfiles}

\begin{document}
% ref: def_3_4
% title: Walks
\def\rowref{def\_3\_4}
\phantomsection\label{def_3_4}

\begin{Def}[Walks]\label{def:walks}
    Let $G = (J, V, E, L)$ be a CDMG (def \ref{def-cdmg}); throughout we use the notation of def \ref{not-cdmg} and the edge classifications of def \ref{def-edge-relations}. Let $v, w \in J \cup V$.

    \begin{enumerate}[label=\roman*.)]
        \item \emph{Walk.}
            A \emph{walk} from $v$ to $w$ in $G$ is, for some integer $n \ge 0$, a tuple
            \[
                (v_0, a_0, v_1, a_1, \dots, v_{n-1}, a_{n-1}, v_n)
            \]
            consisting of $n + 1$ vertices $v_0, v_1, \dots, v_n \in J \cup V$ and $n$ ordered pairs $a_0, a_1, \dots, a_{n-1} \in (J \cup V) \times (J \cup V)$, such that:
            \begin{enumerate}[label=(\alph*)]
                \item $v_0 = v$ and $v_n = w$;
                \item for every $k \in \{0, 1, \dots, n - 1\}$,
                    \[
                        \bigl(a_k = (v_k, v_{k+1}) \;\land\; a_k \in E \cup L\bigr)
                        \;\lor\;
                        \bigl(a_k = (v_{k+1}, v_k) \;\land\; a_k \in E\bigr).
                    \]
            \end{enumerate}
            The case $n = 0$ is admitted: the \emph{trivial walk} is the single-vertex tuple $(v_0)$ with $v_0 = v = w$, no edges, and condition~(b) vacuously satisfied; it is a walk only when $v = w$. For $n \ge 1$, vertices in the tuple may repeat, i.e.\ the same node may appear multiple times among $v_0, \dots, v_n$.

            \emph{End-node classifications.}
            For a non-trivial walk ($n \ge 1$) we further classify behaviour at the two end-nodes:
            \begin{itemize}
                \item the walk is \emph{into $v_0$} iff $a_0$ is an edge into $v_0$ in the sense of def \ref{def-edge-relations}, item~ii.; concretely, iff $\bigl(a_0 = (v_1, v_0) \in E\bigr) \,\lor\, \bigl(a_0 = (v_0, v_1) \in L\bigr)$;
                \item the walk is \emph{out of $v_0$} iff $a_0$ is an edge out of $v_0$ in the sense of def \ref{def-edge-relations}, item~iii.; concretely, iff $a_0 = (v_0, v_1) \in E$;
                \item the walk is \emph{into $v_n$} iff $a_{n-1}$ is an edge into $v_n$ in the sense of def \ref{def-edge-relations}, item~ii.; concretely, iff $\bigl(a_{n-1} = (v_{n-1}, v_n) \in E\bigr) \,\lor\, \bigl(a_{n-1} = (v_{n-1}, v_n) \in L\bigr)$;
                \item the walk is \emph{out of $v_n$} iff $a_{n-1}$ is an edge out of $v_n$ in the sense of def \ref{def-edge-relations}, item~iii.; concretely, iff $a_{n-1} = (v_n, v_{n-1}) \in E$.
            \end{itemize}
            On the trivial walk ($n = 0$) all four classifications are vacuously false (neither $a_0$ nor $a_{n-1}$ exists). On a non-trivial walk with a bidirected first edge ($a_0 = (v_0, v_1) \in L$) the walk is into $v_0$ but not out of $v_0$; hence ``into $v_0$'' and ``out of $v_0$'' are mutually exclusive but not jointly exhaustive over walks (and analogously at $v_n$).

        \item \emph{Directed walk.}
            A \emph{directed walk} from $v$ to $w$ in $G$ is a walk $(v_0, a_0, v_1, \dots, a_{n-1}, v_n)$ from $v$ to $w$ (for some $n \ge 0$) that, in addition, satisfies
            \[
                a_k = (v_k, v_{k+1}) \in E \qquad \text{for every } k \in \{0, 1, \dots, n - 1\}.
            \]
            The trivial walk ($n = 0$, $v = w$) is admitted as a directed walk from $v$ to itself.

        \item \emph{Bidirected walk.}
            A \emph{bidirected walk} from $v$ to $w$ in $G$ is a walk $(v_0, a_0, v_1, \dots, a_{n-1}, v_n)$ from $v$ to $w$ (for some $n \ge 0$) that, in addition, satisfies
            \[
                a_k = (v_k, v_{k+1}) \in L \qquad \text{for every } k \in \{0, 1, \dots, n - 1\}
            \]
            (equivalently, by the symmetry of $L$ from def \ref{def-cdmg}, $(v_{k+1}, v_k) \in L$ for every such $k$). The trivial walk ($n = 0$, $v = w$) is admitted as a bidirected walk from $v$ to itself.

        \item \emph{Collider walk.}
            A \emph{collider walk} from $v$ to $w$ in $G$ is a walk $(v_0, a_0, v_1, \dots, a_{n-1}, v_n)$ from $v$ to $w$ (for some $n \ge 0$) satisfying the case-distinguished constraints below.
            \begin{itemize}
                \item \textbf{Case $n = 0$:} no further constraint; the trivial walk is a collider walk from $v$ to itself when $v = w$.
                \item \textbf{Case $n = 1$:} the single edge $a_0$ is required to be \emph{bidirected}, namely
                    \[
                        a_0 = (v_0, v_1) \in L \qquad \text{(equivalently, } (v, w) \in L\text{).}
                    \]
                    Purely directed edges $(v, w) \in E$ or $(w, v) \in E$ are \emph{not} admitted as collider walks of length~$1$.
                \item \textbf{Case $n \ge 2$:} the edges satisfy:
                    \begin{enumerate}[label=(\alph*)]
                        \item (the first edge places an arrowhead at $v_1$) $\bigl(a_0 = (v_0, v_1) \in E\bigr) \,\lor\, \bigl(a_0 = (v_0, v_1) \in L\bigr)$;
                        \item (every interior edge is bidirected) for every $k \in \{1, 2, \dots, n - 2\}$, $a_k = (v_k, v_{k+1}) \in L$;
                        \item (the last edge places an arrowhead at $v_{n-1}$) $\bigl(a_{n-1} = (v_n, v_{n-1}) \in E\bigr) \,\lor\, \bigl(a_{n-1} = (v_{n-1}, v_n) \in L\bigr)$.
                    \end{enumerate}
                    Under (a)--(c), every interior node $v_k$ (i.e.\ $1 \le k \le n - 1$) has an arrowhead pointing toward it from each of its two incident walk-edges $a_{k-1}$ and $a_k$, recovering the verbal characterisation ``every node strictly between $v$ and $w$ is a collider''.
            \end{itemize}
            \emph{Reconciliation with the source block's $n = 1$ inline note.}
            The source block contains the inline remark ``Note that for $n = 1$ this reads: $v \sus w \in G$''. This remark is \emph{overridden} by the strict reading of Case $n = 1$ above: for $n = 1$, a collider walk requires its unique edge to be bidirected, $a_0 = (v_0, v_1) \in L$. The relaxed reading ``any $v \sus w$ edge'' is \emph{not} adopted here; in particular, neither $(v, w) \in E$ nor $(w, v) \in E$ qualifies as a collider walk of length~$1$.

        \item \emph{Path.}
            A walk $(v_0, a_0, v_1, \dots, a_{n-1}, v_n)$ from $v$ to $w$ in $G$ is called a \emph{path} iff its vertex sequence contains no repetitions, i.e.\ $v_i \ne v_j$ whenever $0 \le i < j \le n$. (Equivalently, the tuple $(v_0, v_1, \dots, v_n)$ consists of $n + 1$ pairwise distinct entries. The trivial walk $n = 0$ is vacuously a path.)

        \item \emph{Bifurcation.}
            A \emph{bifurcation} between $v$ and $w$ in $G$ is a walk $(v_0, a_0, v_1, \dots, a_{n-1}, v_n)$ from $v$ to $w$ for which there exists an index $k$ with $1 \le k \le n$ (in particular $n \ge 1$) such that all of the following hold:
            \begin{enumerate}[label=(\alph*)]
                \item \emph{Distinct end-nodes appearing exactly once.}
                    $v \ne w$, and the end-nodes each occur exactly once in the vertex sequence:
                    \[
                        v_0 \notin \{v_1, v_2, \dots, v_n\}
                        \qquad\text{and}\qquad
                        v_n \notin \{v_0, v_1, \dots, v_{n-1}\}.
                    \]
                \item \emph{Left subwalk is a directed walk from $v_{k-1}$ to $v_0$.}
                    For every $j \in \{0, 1, \dots, k - 2\}$,
                    \[
                        a_j = (v_{j+1}, v_j) \in E.
                    \]
                    (When $k = 1$ the left subwalk is the trivial walk $(v_0)$ and this constraint is vacuous.)
                \item \emph{Right subwalk is a directed walk from $v_k$ to $v_n$.}
                    For every $j \in \{k, k + 1, \dots, n - 1\}$,
                    \[
                        a_j = (v_j, v_{j+1}) \in E.
                    \]
                    (When $k = n$ the right subwalk is the trivial walk $(v_n)$ and this constraint is vacuous.)
                \item \emph{Middle (hinge) edge.}
                    The edge $a_{k-1}$ between $v_{k-1}$ and $v_k$ satisfies
                    \[
                        \bigl(a_{k-1} = (v_k, v_{k-1}) \in E\bigr) \;\lor\; \bigl(a_{k-1} = (v_{k-1}, v_k) \in L\bigr).
                    \]
                \item \emph{(Addition: each end-node has exactly one arrowhead pointing toward it.)}
                    $a_0$ is an edge into $v_0$ and $a_{n-1}$ is an edge into $v_n$, in the sense of def \ref{def-edge-relations}, item~ii.; equivalently, both end-nodes $v_0$ and $v_n$ have exactly one arrowhead pointing toward them (contributed by their unique incident walk-edges $a_0$ and $a_{n-1}$ respectively, uniqueness following from clause~(a)).

                    Spelled out as a case analysis on $k$:
                    \begin{itemize}
                        \item if $k \ge 2$: clause~(b) gives $a_0 = (v_1, v_0) \in E$, which is automatically an edge into $v_0$ (no further restriction);
                        \item if $k = 1$: $a_0 = a_{k-1}$ \emph{is} the middle edge; both alternatives admitted by~(d) — namely $a_0 = (v_1, v_0) \in E$ and $a_0 = (v_0, v_1) \in L$ — are edges into $v_0$, so no further restriction at $v_0$;
                        \item if $k \le n - 1$: clause~(c) gives $a_{n-1} = (v_{n-1}, v_n) \in E$, which is automatically an edge into $v_n$ (no further restriction);
                        \item if $k = n$: $a_{n-1} = a_{k-1}$ \emph{is} the middle edge; among the two alternatives admitted by~(d), only the bidirected alternative $a_{n-1} = (v_{n-1}, v_n) \in L$ is an edge into $v_n$, whereas the directed alternative $a_{n-1} = (v_n, v_{n-1}) \in E$ places its arrowhead at $v_{n-1}$ rather than $v_n$ and is therefore \emph{excluded}.
                    \end{itemize}
                    In particular, the degenerate configuration ``$k = n$ together with the directed hinge $a_{n-1} = (v_n, v_{n-1}) \in E$'' — which would otherwise admit a tuple of the form $v_n \to v_{n-1} \to \cdots \to v_1 \to v_0$ (a plain directed walk from $v_n$ to $v_0$) as a ``bifurcation'' — is \emph{not} a bifurcation under the present definition.
            \end{enumerate}
            \emph{Source of the bifurcation.}
            If clause~(d) is realised by the \emph{directed} alternative $a_{k-1} = (v_k, v_{k-1}) \in E$, then we call $v_k$ the \emph{source} of the bifurcation. If clause~(d) is realised by the \emph{bidirected} alternative $a_{k-1} = (v_{k-1}, v_k) \in L$, no source is defined. Combined with clause~(e), whenever a source $v_k$ is defined we automatically have $1 \le k \le n - 1$ (the case $k = n$ with directed hinge is excluded by~(e)), so the source $v_k$ is an interior vertex of the walk and is distinct from both end-nodes $v_0 = v$ and $v_n = w$.
    \end{enumerate}
\end{Def}

\end{document}
